\documentclass[11pt]{article}

\usepackage[utf8]{inputenc}
\usepackage[T1]{fontenc}
\usepackage{lmodern}
\usepackage[margin=1in]{geometry}
\usepackage{amsmath,amssymb,amsthm}
\usepackage{mathtools}
\usepackage{booktabs}
\usepackage{fancyvrb}
\usepackage{newunicodechar}
\usepackage{xcolor}
\usepackage[hidelinks]{hyperref}
\usepackage{url}
\usepackage{microtype}
\emergencystretch=2em

% Unicode characters that appear in the Lean excerpts.
\newunicodechar{ℝ}{\ensuremath{\mathbb{R}}}
\newunicodechar{ℤ}{\ensuremath{\mathbb{Z}}}
\newunicodechar{≤}{\ensuremath{\le}}
\newunicodechar{⌈}{\ensuremath{\lceil}}
\newunicodechar{⌉}{\ensuremath{\rceil}}
\newunicodechar{ᶜ}{\ensuremath{^{\mathsf c}}}
\newunicodechar{α}{\ensuremath{\alpha}}

\newtheorem{theorem}{Theorem}
\newtheorem{lemma}[theorem]{Lemma}
\newtheorem{corollary}[theorem]{Corollary}
\theoremstyle{definition}
\newtheorem{definition}[theorem]{Definition}
\theoremstyle{remark}
\newtheorem{remark}[theorem]{Remark}

\newcommand{\Gbar}{\overline{G}}
\newcommand{\dbar}{\overline{d}}
\newcommand{\lng}{\operatorname{length}}
\newcommand{\ceil}[1]{\left\lceil #1 \right\rceil}

\fvset{fontsize=\small,frame=single,framesep=6pt,rulecolor=\color{black!40}}

\title{An independent, Lean-verified proof of\\
Written on the Wall II Conjecture 100}
\author{Author: [your name]\\[2pt]
\small Proof search and formalization assisted by Claude (Anthropic)}
\date{October 2026}

\begin{document}
\maketitle

\begin{abstract}
Conjecture~100 of DeLaVi\~na's \emph{Written on the Wall II} (WOWII), a list of
conjectures produced by the program Graffiti.pc, states that every connected
graph $G$ on at least two vertices satisfies
$\alpha(G) \le \bigl\lceil (\max_v l(v) + \tfrac12 \lng(\Gbar))/2 \bigr\rceil$, where
$l(v)$ is the independence number of the neighbourhood of $v$ and
$\lng(\Gbar)$ is the Euclidean norm of the degree sequence of the complement.
\textbf{This is not the first proof.} The conjecture was proved by Kias
Henry~\cite{henry} (Zenodo, August 2026; the proof was reviewed by
E.~DeLaVi\~na, who marked the conjecture as true), and a Lean formalization by
the GitHub user vatsj was submitted to the \texttt{formal-conjectures}
repository as pull request \#5221~\cite{vatsj}. We give an independent
argument that does not use the Cauchy--Schwarz or AM--GM inequalities: after
elementary bounds on complement degrees and a double count, everything reduces
to the positivity of one explicit bivariate polynomial
$P(k,t)$ on the positive integers. We also describe a self-contained Lean~4 /
Mathlib formalization of this argument (about 220 lines), checked against
the unchanged statement in \texttt{formal-conjectures}.
\end{abstract}

\section{Introduction}

Graffiti.pc is a conjecture-generating program written by Ermelinda
DeLaVi\~na; its output is collected in the online list \emph{Written on the
Wall II} (WOWII)~\cite{wowii}. Google DeepMind's \path{formal-conjectures}
repository~\cite{fc} contains Lean~4 statements of many WOWII conjectures,
each tagged with a status. Conjecture~100 concerns the independence number,
local independence numbers, and the degree sequence of the complement graph.

\paragraph{Prior work.}
We state plainly that the result in this note is \emph{already known}:
\begin{itemize}
\item Kias Henry, \emph{A Proof of Written on the Wall II Conjecture 100},
  Zenodo, 13~August~2026, \href{https://doi.org/10.5281/zenodo.21914031}{doi:10.5281/zenodo.21914031}~\cite{henry}.
  Henry's proof uses the Cauchy--Schwarz and AM--GM inequalities. It was
  reviewed by E.~DeLaVi\~na, who marked the conjecture as true on the WOWII
  page.
\item A Lean proof by the GitHub user vatsj, submitted as pull
  request~\#5221 to the \path{formal-conjectures} repository~\cite{vatsj}
  (opened 31~August~2026; not merged as of 9~October~2026).
\end{itemize}
At the commit of \texttt{formal-conjectures} we worked against, the statement
was still tagged \texttt{research open}; that is why we took it up, and we
learned of the two works above afterwards. Our argument was found
independently of both.

\paragraph{What this note contributes.}
Given the above, our contribution is modest:
\begin{enumerate}
\item A second, independent proof. It uses only pointwise lower bounds on
  complement degrees, the inequalities $(x+y)^2 \ge x^2 + 2xy$ and
  $c^2 \ge tc$ (for $c \ge t \ge 0$), and one double count. No
  Cauchy--Schwarz or AM--GM step is needed. All of the arithmetic is collected
  into a single polynomial inequality $P(k,t) > 0$ (Lemma~\ref{lem:poly}).
\item A compact formalization in Lean~4 with Mathlib (about 220 non-blank
  lines including the three arithmetic lemmas), which fills the
  \texttt{sorry} in the existing \texttt{formal-conjectures} file without
  changing the statement.
\end{enumerate}
We have not compared our formalization with the one in~\cite{vatsj} line by
line, and we make no claim about which is preferable.

\section{Statement}

All graphs are finite and simple. For a graph $G$ with vertex set $V$ and $n = |V|$ we
write $N(v)$ for the (open) neighbourhood of $v$, $\deg v = |N(v)|$,
$\alpha(G)$ for the independence number, and $\Gbar$ for the complement, so
that $\deg_{\Gbar}(v) = n-1-\deg v$.

\begin{definition}
For $v \in V$, let $l(v) = \alpha\bigl(G[N(v)]\bigr)$ be the independence
number of the subgraph induced by the neighbourhood of $v$ (with $l(v)=0$
if $N(v)=\emptyset$). The \emph{length} of a graph $H$ is the Euclidean norm
of its degree sequence,
$\lng(H) = \bigl(\sum_{v} \deg_H(v)^2\bigr)^{1/2}$.
\end{definition}

\begin{theorem}[WOWII Conjecture 100]\label{thm:main}
Let $G$ be a connected graph with at least two vertices. Then
\[
  \alpha(G) \;\le\; \ceil{\frac{\max_{v} l(v) + \tfrac12\,\lng(\Gbar)}{2}}.
\]
\end{theorem}

\begin{remark}
The hypothesis $n \ge 2$ is needed. For the one-vertex graph $K_1$ we have
$\alpha = 1$, $l(v) = 0$ and $\lng(\Gbar) = 0$, so the right-hand side is
$0$. The WOWII text says only ``simple connected graph''; the
\texttt{formal-conjectures} statement assumes the vertex type is nontrivial,
and we do the same.
\end{remark}

\section{The proof}

Fix $G$ as in Theorem~\ref{thm:main}. Throughout we use the following
notation.
\begin{itemize}
\item $I$ is a maximum independent set, $a = |I| = \alpha(G)$.
\item $O = V \setminus I$ and $m = |O|$.
\item $K = \max_v l(v)$.
\item $y(w) = |N(w) \cap I|$ for $w \in O$.
\item $E = $ the number of edges between $I$ and $O$.
\item $S = \sum_{v \in V} \deg_{\Gbar}(v)^2$, so $\lng(\Gbar) = \sqrt{S}$.
\item $D = 2a - 2 - K$.
\end{itemize}

\subsection{Reduction}

Since $a$ is an integer, $\ceil{x} \ge a$ holds if and only if $x > a-1$.
With $x = (K + \tfrac12\sqrt S)/2$ this is equivalent to
\begin{equation}\label{eq:goal}
  \sqrt{S} \;>\; 2D = 2(2a-2-K).
\end{equation}
If $D < 0$ then \eqref{eq:goal} is immediate because $\sqrt S \ge 0$. So
from now on we assume $D \ge 0$, and it suffices to prove
\begin{equation}\label{eq:goal2}
  S \;>\; 4D^2.
\end{equation}

\subsection{Elementary facts}

\begin{lemma}\label{lem:facts}
With the notation above:
\begin{enumerate}
\item[(1)] $y(w) \le l(w) \le K$ for every $w \in O$.
\item[(2)] $a \le E \le K m$.
\item[(3)] For $u \in I$, $\deg_{\Gbar}(u) \ge (a-1) + (m - \deg u)$, and
  $m - \deg u \ge 0$. For $w \in O$, $\deg_{\Gbar}(w) \ge a - y(w)$.
\end{enumerate}
\end{lemma}

\begin{proof}
(1) The set $N(w) \cap I$ is independent in $G$ (it is a subset of $I$) and is
contained in $N(w)$, so it is an independent set of $G[N(w)]$. Hence
$y(w) \le \alpha(G[N(w)]) = l(w) \le K$.

(2) Because $I$ is independent, every neighbour of a vertex $u \in I$ lies in
$O$; hence $\deg u = |N(u) \cap O|$ and $E = \sum_{u \in I} \deg u$. Since $G$ is
connected and has at least two vertices, every vertex has degree at least~$1$,
so $E \ge |I| = a$. Counting the same edges from the $O$ side gives
$E = \sum_{w \in O} y(w) \le K m$ by~(1).

(3) For $u \in I$, the vertices of $I \setminus \{u\}$ are not adjacent to $u$
in $G$, and neither are the vertices of $O \setminus N(u)$. Since
$N(u) \subseteq O$, the latter set has $m - \deg u \ge 0$ elements. (In fact
equality holds, but we only need the inequality.) For $w \in O$, the
$a - y(w)$ vertices of $I \setminus N(w)$ are non-neighbours of $w$ distinct
from $w$.
\end{proof}

Summing $m - \deg u$ over $u \in I$ and using (2) gives
\begin{equation}\label{eq:amE}
  \sum_{u \in I} (m - \deg u) \;=\; am - E \;\ge\; 0.
\end{equation}

\subsection{Case \texorpdfstring{$K \ge a$}{K >= a}}

Here $0 \le D = 2a - 2 - K \le a - 2$, so $a \ge 2$. Keeping only the vertices
of $I$ and using Lemma~\ref{lem:facts}(3),
\[
  S \;\ge\; \sum_{u \in I} \deg_{\Gbar}(u)^2 \;\ge\; a(a-1)^2.
\]
It remains to check $a(a-1)^2 > 4(a-2)^2 \ge 4D^2$. For $a = 2$ this reads
$2 > 0$; for $a = 3$ it reads $12 > 4$; and for $a \ge 4$ we have
$a(a-1)^2 \ge 4(a-1)^2 > 4(a-2)^2$. This proves~\eqref{eq:goal2}.

\subsection{Case \texorpdfstring{$K < a$}{K < a}}

Put $t = a - K \ge 1$. Note $K \ge 1$: otherwise $E \le Km = 0 < a$,
contradicting Lemma~\ref{lem:facts}(2).

\emph{Contribution of $I$.} For $x, y \ge 0$ we have
$(x+y)^2 \ge x^2 + 2xy$. Applying this with $x = a-1$ and
$y = m - \deg u$ (Lemma~\ref{lem:facts}(3)) and summing with~\eqref{eq:amE},
\[
  \sum_{u \in I} \deg_{\Gbar}(u)^2
  \;\ge\; a(a-1)^2 + 2(a-1)(am - E).
\]

\emph{Contribution of $O$.} For $w \in O$ write $c_w = \deg_{\Gbar}(w)$. By
Lemma~\ref{lem:facts}(1),(3), $c_w \ge a - y(w) \ge a - K = t > 0$, hence
$c_w^2 \ge t\,c_w \ge t\,(a - y(w))$. Summing over $O$ and using
$\sum_{w \in O} y(w) = E$,
\[
  \sum_{w \in O} \deg_{\Gbar}(w)^2 \;\ge\; t\,(am - E).
\]

\emph{Combining.} Adding the two bounds,
\[
  S \;\ge\; a(a-1)^2 + (2a - 2 + t)(am - E).
\]
By Lemma~\ref{lem:facts}(2), $am - E \ge am - Km = tm$, and
$2a-2+t \ge 0$, so
\begin{equation}\label{eq:S}
  S \;\ge\; a(a-1)^2 + (2a-2+t)\,t\,m.
\end{equation}
Multiply~\eqref{eq:S} by $K > 0$ and use $Km \ge a$ (again Lemma~\ref{lem:facts}(2)):
\[
  K S \;\ge\; K a (a-1)^2 + (2a-2+t)\,t\,a.
\]
Since $D = 2a - 2 - K = a - 2 + t$, we get
\begin{equation}\label{eq:KS}
  K S - 4 K D^2 \;\ge\; K a (a-1)^2 + (2a - 2 + t)\,a\,t - 4K(a - 2 + t)^2 .
\end{equation}
Substituting $a = K + t$ and expanding, the right-hand side
of~\eqref{eq:KS} equals $P(K,t)$, where
\begin{equation}\label{eq:P}
\begin{aligned}
  P(k,t) \;=\;& k\,(k^3 - 6k^2 + 17k - 16)
  \;+\; t\,k\,(3k^2 - 18k + 31) \\
  &+\; t^2\,(3k^2 - 13k - 2) \;+\; t^3\,(k+3).
\end{aligned}
\end{equation}
(We checked this identity by symbolic expansion in SymPy. In the Lean
formalization the lemma \texttt{key\_poly} is stated in the unexpanded form of
the right-hand side of~\eqref{eq:KS}, so the expansion is not a separate
trusted step there.) By
Lemma~\ref{lem:poly} below, $P(K,t) > 0$, so $K S > 4 K D^2$, and dividing
by $K > 0$ gives~\eqref{eq:goal2}.

\subsection{The polynomial inequality}

\begin{lemma}\label{lem:poly}
$P(k,t) > 0$ for all integers $k \ge 1$ and $t \ge 1$.
\end{lemma}

\begin{proof}
\emph{The case $k \ge 5$.} Every term of~\eqref{eq:P} is then positive:
\begin{itemize}
\item $k^3 - 6k^2 + 17k - 16$ has a single real root, near $1.61$, and is
  positive for $k \ge 2$ (its values at $k = 2,3,4,5$ are $2, 8, 20, 44$);
\item $3k^2 - 18k + 31$ has discriminant $324 - 372 = -48 < 0$, so it is
  positive for all real $k$;
\item $3k^2 - 13k - 2$ has its larger root near $4.48$, so it is positive for
  $k \ge 5$ (its value at $k=5$ is $8$);
\item $k + 3 > 0$.
\end{itemize}

\emph{The cases $k \le 4$.} Here $P(k,\cdot)$ is an explicit cubic. Writing
$t = s + 1$ with $s \ge 0$:
\begin{center}
\begin{tabular}{ccc}
\toprule
$k$ & $P(k,t)$ & in terms of $s = t-1 \ge 0$ \\
\midrule
$1$ & $4t^3 - 12t^2 + 16t - 4$ & $4s^3 + 4s + 4$ \\
$2$ & $5t^3 - 16t^2 + 14t + 4$ & $5s^3 - s^2 - 3s + 7$ \\
$3$ & $6t^3 - 14t^2 + 12t + 24$ & $6s^3 + 4s^2 + 2s + 28$ \\
$4$ & $7t^3 - 6t^2 + 28t + 80$ & $7s^3 + 15s^2 + 37s + 109$ \\
\bottomrule
\end{tabular}
\end{center}
For $k = 1, 3, 4$ all coefficients in $s$ are positive, so $P > 0$. For
$k = 2$: at $s = 0$ the value is $7$, and for integers $s \ge 1$ we have
$5s^3 - s^2 - 3s = s(5s^2 - s - 3) \ge 1$, so the value is at least $8$.
\end{proof}

\begin{remark}
The arithmetic margin is small. By the case analysis above, the minimum of
$P$ over positive integers is $P(1,1) = 4$ (the next smallest values with
$k \le 4$ are $P(2,1) = 7$ and $P(2,2) = 8$, and $P(k,t) \ge 5 \cdot 44$ for
$k \ge 5$). This says nothing about the tightness of Theorem~\ref{thm:main}
itself, since the estimates leading to~\eqref{eq:S} and~\eqref{eq:KS} lose
a lot for most graphs.
\end{remark}

\subsection{Conclusion}

In both cases~\eqref{eq:goal2} holds whenever $D \ge 0$, hence
$\sqrt{S} > 2(2a - 2 - K)$, i.e.
$(K + \tfrac12\sqrt S)/2 > a - 1$, and therefore
$\ceil{(K + \tfrac12 \lng(\Gbar))/2} \ge a = \alpha(G)$. This proves
Theorem~\ref{thm:main}. \qed

\section{Lean formalization}

\subsection{Setting}

The formalization lives in the file
\begin{center}\small\path{FormalConjectures/WrittenOnTheWallII/GraphConjecture100.lean}\end{center}
of \path{google-deepmind/formal-conjectures}~\cite{fc}, at commit
\begin{center}\small\texttt{b3f26411f06b0b0ec1923da80b62585f23288d39},\end{center}
built with Lean~4
v4.33.1~\cite{lean4} and the Mathlib~\cite{mathlib} version pinned by that
commit. The theorem statement is the repository's own, unchanged; we only
replaced its \texttt{sorry} by a proof and changed the category attribute
from \texttt{research open} to \texttt{research solved}. The statement is:

\begin{Verbatim}
theorem conjecture100 (G : SimpleGraph α) [DecidableRel G.Adj]
    (h : G.Connected) :
    let maxL := (Finset.univ.image (indepNeighborsCard G)).max' (by simp)
    (G.indepNum : ℝ) ≤
      ⌈((maxL : ℝ) + (1 / 2) * (degreeL2Norm Gᶜ : ℝ)) / 2⌉
\end{Verbatim}

\noindent
Here the vertex type \texttt{α} carries \texttt{[Fintype α] [DecidableEq α]
[Nontrivial α]}, so $G$ has at least two vertices.
\texttt{indepNeighborsCard G v} is the independence number of the
neighbourhood of \texttt{v}, and \texttt{degreeL2Norm H} is
$\sqrt{\sum_v \deg_H(v)^2}$. Both are defined in the repository's shared
utilities. The command \texttt{\#print axioms conjecture100} reports only
\texttt{propext}, \texttt{Classical.choice} and \texttt{Quot.sound}, the
standard axioms of Lean and Mathlib; in particular there is no
\texttt{sorryAx}.

\subsection{Structure}

The proof follows Section~3 closely. Three private lemmas over $\mathbb{Z}$
carry the arithmetic, and the main theorem does the combinatorics.

\begin{itemize}
\item \texttt{key\_poly (k t : ℤ)}: for $k, t \ge 1$,
  $4k(k + 2t - 2)^2 < k(k+t)(k+t-1)^2 + (2(k+t) - 2 + t)(k+t)t$, which is
  $P(k,t) > 0$ in unexpanded form. The proof splits on $k \le 4$ versus
  $k \ge 5$. For $k \le 4$, \texttt{interval\_cases k} produces the four
  cubics, each closed by \texttt{nlinarith} with nonnegative products of
  $t-1$ as hints. For $k \ge 5$, \texttt{nlinarith} is given products of
  $k-5 \ge 0$ and $t \ge 0$ as hints. (So the Lean certificate for $k \ge 5$
  is a Positivstellensatz-style combination found by \texttt{nlinarith}, not
  literally the term-by-term argument of Lemma~\ref{lem:poly}.)
\item \texttt{arith\_lt}: the case $K < a$. From the integer hypotheses
  $a \le E \le Km$, the $I$-bound
  $a(a-1)^2 + 2(a-1)(am-E) \le X$ and the $O$-bound $(a-K)(am-E) \le Y$, it
  derives $4(2a-2-K)^2 < X + Y$, following~\eqref{eq:S}--\eqref{eq:KS} and
  calling \texttt{key\_poly}.
\item \texttt{arith\_ge}: the case $a \le K$. From $a(a-1)^2 \le X$ and
  $0 \le 2a-2-K$ it derives $4(2a-2-K)^2 < X$.
\item \texttt{conjecture100}: obtains $I$, defines $y$, proves
  Lemma~\ref{lem:facts} and the two sums, splits into the three cases
  ($D < 0$, $K < a$, $K \ge a$), and finally passes from the integer
  inequality $4D^2 < S$ to the real inequality with the ceiling.
\end{itemize}

\subsection{Mathlib API used}

The combinatorial steps reduce to a handful of Mathlib lemmas.
\begin{center}
\small
\begin{tabular}{p{0.5\textwidth}p{0.38\textwidth}}
\toprule
Mathlib declaration & Used for \\
\midrule
\texttt{SimpleGraph.exists\_isNIndepSet\_indepNum} & a maximum independent set $I$, $|I| = a$ \\
\texttt{IsIndepSet.card\_le\_indepNum} & $y(w) \le l(w)$ and $a \ge 1$ \\
\texttt{Preconnected.degree\_pos\_of\_nontrivial} & every degree is $\ge 1$ \\
\texttt{neighborFinset\_compl} & neighbours in $\Gbar$ as $N(v)^{\mathsf c} \setminus \{v\}$ \\
\texttt{Finset.sum\_card\_bipartiteAbove\_eq\_\allowbreak sum\_card\_bipartiteBelow} & double counting $E$ \\
\texttt{Int.le\_ceil\_iff}, \texttt{Real.lt\_sqrt\_of\_sq\_lt} & the final reduction~\eqref{eq:goal} \\
\texttt{nlinarith}, \texttt{interval\_cases}, \texttt{omega} & arithmetic \\
\bottomrule
\end{tabular}
\end{center}

\noindent
The whole proof, including the three arithmetic lemmas, occupies about 220
non-blank lines. To check it:
\begin{Verbatim}
git clone https://github.com/google-deepmind/formal-conjectures
cd formal-conjectures
git checkout b3f26411f06b0b0ec1923da80b62585f23288d39
git apply ../formal-conjectures.patch
lake exe cache get
lake build FormalConjectures.WrittenOnTheWallII.GraphConjecture100
\end{Verbatim}

\section{Computational check}

Before the proof was found, the inequality was tested exhaustively on all
connected graphs with $2 \le n \le 9$ vertices, generated with
\texttt{geng} from nauty~\cite{nauty}. There are
$1 + 2 + 6 + 21 + 112 + 853 + 11117 + 261080 = 273{,}192$ such graphs up to
isomorphism. No counterexample was found. For this note we repeated the check
with a separate script that avoids floating point: by the reduction
in~\S3.1, the inequality holds if and only if $4a - 4 - 2K < 0$ or
$S > (4a-4-2K)^2$, which is an integer comparison. The repeat also found no
counterexample. This is evidence, not part of the proof; Graffiti.pc
conjectures are by construction consistent with the program's own database
of graphs.

\section{Discussion}

We have not studied the argument of~\cite{henry} in detail, beyond noting
that it uses the Cauchy--Schwarz and AM--GM inequalities, so we do not
attempt a detailed comparison. Our proof uses only the termwise bounds of
Lemma~\ref{lem:facts}(3), the two elementary inequalities
$(x+y)^2 \ge x^2 + 2xy$ and $c^2 \ge tc$, and the double count
$a \le E \le Km$; the price is the case analysis in Lemma~\ref{lem:poly}.
One practical consequence of this shape is that the formal proof needs no
real analysis until the last step: every inequality before it is between
integers, which \texttt{nlinarith} and \texttt{omega} handle directly.

\paragraph{Acknowledgements.}
The proof search and the Lean formalization were carried out with the
assistance of Claude (Anthropic). We thank the maintainers of
\texttt{formal-conjectures} for the formal statement and E.~DeLaVi\~na for
maintaining WOWII.

\begin{thebibliography}{9}

\bibitem{wowii}
E.~DeLaVi\~na,
\emph{Written on the Wall II (Conjectures of Graffiti.pc)},
\url{http://cms.dt.uh.edu/faculty/delavinae/research/wowII/}.

\bibitem{henry}
K.~Henry,
\emph{A Proof of Written on the Wall II Conjecture 100},
Zenodo, 13 August 2026,
\href{https://doi.org/10.5281/zenodo.21914031}{doi:10.5281/zenodo.21914031}.

\bibitem{vatsj}
vatsj,
\emph{Lean proof of WOWII Conjecture 100},
pull request \#5221 to \path{google-deepmind/formal-conjectures},
opened 31 August 2026,
\url{https://github.com/google-deepmind/formal-conjectures/pull/5221}.

\bibitem{fc}
Google DeepMind,
\emph{formal-conjectures}: a collection of formalized conjectures in Lean,
\url{https://github.com/google-deepmind/formal-conjectures}.

\bibitem{lean4}
L.~de~Moura and S.~Ullrich,
The Lean~4 theorem prover and programming language,
in \emph{Automated Deduction -- CADE 28}, Lecture Notes in Comput. Sci.
12699, Springer, 2021, 625--635.

\bibitem{mathlib}
The mathlib Community,
The Lean mathematical library,
in \emph{Proceedings of the 9th ACM SIGPLAN International Conference on
Certified Programs and Proofs (CPP 2020)}, ACM, 2020, 367--381.

\bibitem{nauty}
B.~D. McKay and A.~Piperno,
Practical graph isomorphism, II,
\emph{J. Symbolic Comput.} \textbf{60} (2014), 94--112.

\end{thebibliography}

\end{document}
